
When LLM-generated code fails, developers can provide feedback from static analysis tools (pylint, mypy) or from test execution. Prior work demonstrates that both feedback types improve code quality, but whether they detect the same bugs or different ones remains unquantified. This study provides the first controlled measurement of overlap between static-detected and execution-detected error sets. Analyzing 542 problems from HumanEval+ and MBPP+, we observe a Jaccard similarity of 0.0 between the two error classes---complete orthogonality with no overlap. Static analysis detects structural errors (types, undefined variables, style violations) in 98.6\% of test-failing code. However, the complementary hypothesis---that execution detects behavioral errors in static-clean code---was not validated: only 0.3\% of static-clean canonical solutions failed execution tests, far below the 40\% threshold. This discrepancy arises because canonical solutions are designed to pass tests, not because execution feedback lacks value. The complete separation observed between error classes supports the theoretical basis for combining feedback types in iterative refinement systems, though empirical validation on actual LLM-generated code remains necessary.

